\chapter{系统架构与工程实现}

\section{总体架构}

图 \ref{fig:system-architecture} 展示系统架构。学生在自己的仓库中维护规格、实现代码、测试和课程运行配置。\code{vos-cli} 提供本地命令入口，\code{vos-core} 负责规格检查、Agent 调度、运行编排、证据记录、报告和提交，\code{vos-spec} 解析五类规格，\code{vos-runtime} 执行宿主、QEMU 与开发板任务，\code{vos-kb} 管理课程知识来源。\code{vos-agent} 在进程内运行角色化模型会话。Portal 与独立静态 Demo 作为已有应用继续维护和验证。

\begin{figure}[htbp]
  \centering
  \includegraphics[width=\textwidth]{system-architecture.pdf}
  \caption{VeriSpecOSLab 当前系统架构与学生主链}
  \label{fig:system-architecture}
\end{figure}

\section{学生主链的状态转换}

\subsection{初始化与设计}

开始一个 Lab 时，学生先用 \code{vos init} 创建实验所需的基本目录和配置，再通过 \code{agent ask} 讨论课程概念和设计选择。学生随后亲手写下系统采用的总体方案，以及当前模块负责什么、与哪些模块交互、准备怎样验收。平台检查这些说明是否完整一致，评审助手再指出需要补充或澄清的内容。学生理解并修改后，将规格提交到 Git，后续代码实现便以这份设计为依据。

\subsection{规格与实现}

学生完成设计说明并提交后，可以让实现助手处理当前模块。平台先从规格中确定本次 Lab 涉及的代码范围，再在与学生工作目录分开的临时目录中生成代码和配套测试。生成结果需要通过编译和多类测试，平台还会核对实际修改是否超出当前任务。代码范围和测试结果符合要求后，这次实现会作为一份独立提交加入学生仓库。学生由此得到一份与自己设计对应并经过实际验证的代码修改。

\subsection{运行、验证与提交}

开发过程中，学生可以随时编译系统、启动 QEMU 或连接开发板，根据运行现象继续调试。准备正式验收时，学生先把当前代码提交到 Git，平台再运行完整验证，并将代码版本、规格、实际命令、日志和测试结果对应起来。\code{vos report} 汇总这些实验记录，\code{vos submit} 打包同一版本的代码和材料。教师拿到提交后，可以从报告结论回到相应代码与运行过程进行复查。

\section{技术栈与代码规模}

平台使用 TypeScript 和 Bun 实现。排除 \code{node\_modules} 后，\code{vos/} 下共有 338 个 TypeScript 文件、约 71,401 行，其中 127 个测试文件约 22,283 行，覆盖 CLI、Agent、共享包、Portal、Demo 和测试。

\section{CLI 入口与命令执行}

\code{vos-cli} 负责可执行入口和包契约。命令解析和类型化执行位于 \code{vos-core}，测试可以直接构造命令对象。课程部署从 \code{vos/apps/vos-cli} 执行 \code{bun link}，链接后的命令始终运行当前仓库中的代码，使开发、教学和复现使用同一版本。

公开命令覆盖初始化与诊断、规格、Agent、知识库、构建、运行、验证、报告和提交。内部仍保留兼容类型与 Portal 命令解析测试，学生实际使用的命令以 \code{vos --help} 为准。

\section{错误模型与进度事件}

核心命令返回 \code{CommandOutcome} 和明确状态。验证失败、策略阻止执行、运行超时、用户取消和内部错误会在上层保持区分。进度界面根据 TTY、JSON、CI 和 \code{--progress} 选项决定是否渲染动态行。相同的底层进度也可以通过 MCP \code{report\_progress} 由 Agent 上报。

结构化错误便于报告按真实原因汇总结果。例如，``工作树存在未提交改动''记为 \code{policy\_blocked}，不计作测试失败。``模型没有提交符合协议的结构化结果''属于 Agent 协议错误，也不计作内核构建失败。

\section{配置解析与接口检查}

规格、运行配置、Agent 配置、结果 schema 和 HTTP contract 默认拒绝未知字段。字段发生变化时，schema、测试和文档需要同步更新，拼写错误或版本差异会在解析阶段直接显示。\code{vos.yaml} 的版本结构、知识库来源、目标路径和 Agent 配置字段均有对应测试。

配置中的凭据通过环境变量名引用，TOML 解析器负责检查字段和格式。Portal 与 Server 为每条 API 路由建立类型化 contract，并检查访问凭据和产物路径。

\section{测试与构建组织}

测试使用 Bun 内置 runner，并放在所属模块的测试目录中。测试覆盖以下层次：

\begin{itemize}
  \item schema 与 parser 单元测试，检查严格字段、YAML、TOML、路径和引用。
  \item 流程约束测试，检查 clean HEAD、文件修改范围、SpecPatch、凭据与测试可见性。
  \item 学生流程集成测试，覆盖初始化、手写规格检查与评审、临时实现目录、测试目标更新、构建、QEMU、验证、报告和提交。
  \item Agent 协议测试，覆盖 Provider 配置、工具 profile、MCP 提交、循环上限和 TUI。
  \item Portal 测试，覆盖接口契约、存储 adapter、鉴权和外部服务集成接口。
  \item 课程历史审计，逐个标签检查提前出现的路径、ID、测试和术语。
\end{itemize}

PostgreSQL、Gitea、MinIO、OIDC、Docker 和真实模型分别具有独立的集成测试入口。fixture 与 fake runner 用于快速检查控制流程，配置相应服务后可以继续运行连通与端到端测试。

workspace 的 \code{bun run build} 先检查应用引用关系，再分别构建 CLI、Agent、Demo 与 Portal。CLI 和 Agent 使用 Bun compile 生成本地可执行文件，Demo 和 Portal 使用 Vite 与 TypeScript。所有应用共同参加类型检查和构建。

